\section{The reference unipotent covolume and its exact discriminant}
\label{sec:reference-covolume-v56}

This chapter proves the original \texttt{CovolRef}.  The calculation concerns
the prescribed reference measure obtained from the Euclidean volume of the
standard nilpotent Lie algebra.  It is a separate calculation from the
adelic volume normalization and the actual \texttt{InVolumes} established
in GMR and SVW.  In particular, none of the supplied measures
\(X.\mu_N\), \(X.\mu_H\), the global character, or the adelic bridge is
changed here.

Fix any original \(D7\) datum \(D\) and any supplied
\(\mathrm{HidaSetting}(D)\), denoted by \(X\).  We work on its original
Hermitian space \((V,h)\), with \(h\) conjugate-linear in its first
variable.  Write
\[
 \begin{gathered}
 E=D.host.E,\quad F=D.host.F,\quad r=D.R.rS,\quad k=D.R.k,\\
 q=k+1,\quad d=[F:\mathbb Q],\quad
 \Delta=|\operatorname{disc}(X.B.I.F)|.
 \end{gathered}
\]
The supplied coefficient isomorphism identifies \(X.B.I.F\) with \(F\),
so this is also \(|\operatorname{disc}F|\).  We retain the supplied
ordering of the \(d\) real places and its chosen complex embeddings
\(\sigma_v:E\hookrightarrow\mathbb C\).  Put
\[
 \mathfrak d=\mathfrak D_{F/\mathbb Q}\mathcal O_E,
 \qquad \mathfrak a_b=\mathfrak d^{-(r-b)}\quad(0\leq b<r),
 \qquad T_w=\mathrm{Tw}(w)>0.
\tag{RCV.1}
\]
These are the literal \texttt{ddS}, \texttt{aaD}, and \texttt{Tw} of the
source.  Its hyperbolic vectors will be written \(e_b,f_b\), its
orthogonal-complement lattice \(L_b\), and its full lattice \(L_0\).
Thus the supplied lattice identity is
\[
 L_0=\bigoplus_{b<r}\mathfrak a_b e_b\ \oplus L_b\ \oplus
       \bigoplus_{b<r}\mathfrak a_b^{-1}f_b,
 \qquad h(e_b,f_a)=\delta_{ba}.
\tag{RCV.2}
\]
The directness in this display will also be verified below.  The letter
\(b\) in \(L_b\) is part of the source's name for one lattice; it does
not index a family of lattices.

For each word \(w\), the supplied vectors
\((w_{w,v,1},\ldots,w_{w,v,k},z_w)\) form a basis of
\(W_v=\langle e,f\rangle^\perp_v\), with Gram matrix
\(\operatorname{diag}(1,\ldots,1,-T_w)\).  We shall use two words
independently: \(u\) specifies a reference measure, and \(w\) specifies
the coordinate box and the arithmetic fundamental domain.  No equality
between them is assumed.

All source statements are read at commit
\texttt{7cbf7a0c32fca1badf05931879a37c3f8a2b03b4}.  The accompanying
manifest gives exact declaration ranges.  The only inherited human
results used below are ASG2, for the original trace-one element and
unramified CM extension, and AMD8, for Haar uniqueness at the source's
weak \texttt{IsHaarOnSet} scope.  The intervening arithmetic and
coordinate arguments are proved in this chapter.

\subsection{The four arithmetic coordinate lattices}

% BEGIN RCV1
\paragraph{RCV1. Discriminants, extended ideals, and conjugation.}
\label{claim:RCV1}
For the actual coefficient fields and every \(b<r\),
\[
 |\operatorname{disc}E|=\Delta^2,\qquad
 N_E(\mathfrak d)=\Delta^2,\qquad
 N_E(\mathfrak a_b)=\Delta^{-2(r-b)}.
\tag{RCV.3}
\]
All these ideals are nonzero and conjugation invariant.  There is an
integral \(\theta\in E\) with \(\theta+\bar\theta=1\).
Here \(N_E\) denotes the positive absolute norm of a fractional ideal.
\begin{proof}
ASG2 applies to the original host and gives an integral trace-one
element.  The source's \(\mathrm{theta}(D)\) chooses an element with
exactly these two properties when one exists; we take that actual
choice as \(\theta\).  ASG2 also gives conjugation invariance and
absence of finite ramification in
\(E/F\).  We recall the determinant consequence to specify the exact
normalization being used.  For a finite extension of number fields,
the discriminant of the integral trace pairing satisfies
\[
 |\operatorname{disc}E|
 =|\operatorname{disc}F|^{[E:F]}
       N_F(\mathfrak D_{E/F}^{\mathrm{disc}}).
\]
This follows by writing the trace as the composite of the two traces:
in local integral bases the determinant is the determinant from the
base field repeated \([E:F]\) times, multiplied by the norm of the
relative trace determinant.  Changing a local basis multiplies the
determinant by the square of a unit, so the local computations glue
as fractional ideals.  At a split or unramified quadratic place the
integral relative trace pairing is perfect.  Its determinant is a
unit, including at residue characteristic two.  Consequently the
relative discriminant ideal is the unit ideal, and \([E:F]=2\) gives
the first equality in (RCV.3).

For a nonzero ideal \(J\) of \(\mathcal O_F\),
\(N_E(J\mathcal O_E)=N_F(J)^2\).  To see this without a global
free-basis assumption, localize at each prime of \(\mathcal O_F\).
The integral algebra \(\mathcal O_E\) is a free module of rank two
there, and its quotient by a power of the prime has twice the
module length of the corresponding quotient of the base ring.
Multiplying residue-field cardinalities proves the norm identity.
It extends to fractional ideals by taking quotients.
The trace-dual definition of the different gives
\(N_F(\mathfrak D_{F/\mathbb Q})=|\operatorname{disc}F|\): in a
\(\mathbb Z\)-basis, the index of the trace dual is the absolute
determinant of the trace Gram matrix.  Apply the extended-ideal norm
identity and then take the prescribed negative powers.  This proves
the remaining equalities.  Finally, the supplied field isomorphism
preserves integral elements and the rational trace pairing; it
therefore preserves its discriminant.  Thus every occurrence of
\(\Delta\) in (RCV.3) is the discriminant of the field in the actual
consumer, not of a replacement field.
\end{proof}
% END RCV1

% BEGIN RCV2
\paragraph{RCV2. Euclidean covolumes with every measurable domain.}
\label{claim:RCV2}
Let \(\Lambda\) be the integer span of a real basis
\(v_1,\ldots,v_m\) of a finite-dimensional real space with its specified
Lebesgue measure.  It is countable and discrete, has a Borel strict
additive fundamental domain, and every measurable strict additive
fundamental domain has volume \(|\det(v_1,\ldots,v_m)|\).
In the chosen embeddings of a CM field this gives, for every nonzero
fractional ideal \(I\subset E\),
\[
 \operatorname{covol}(\sigma(I))
 =2^{-d}\sqrt{|\operatorname{disc}E|}\,N_E(I).
\tag{RCV.4}
\]
For a fractional ideal \(J\subset F\), its real embedding lattice has
covolume \(\sqrt{\Delta}\,N_F(J)\).
\begin{proof}
The image of \([0,1)^m\) under the basis matrix is a Borel fundamental
parallelotope: coordinatewise integer subtraction gives existence
and the half-open intervals give uniqueness.  The same invertible
matrix proves discreteness and the determinant formula.  If \(A\)
is any other measurable strict fundamental domain and \(P\) is this
parallelotope, partition \(A\) by the countable sets
\(A\cap(\lambda+P)\), \(\lambda\in\Lambda\).  Translating these
pieces by \(-\lambda\) partitions \(P\).  Translation invariance
and countable additivity give \(\operatorname{vol}A=
\operatorname{vol}P\).  This argument first takes place in
\([0,\infty]\); it does not assume that \(A\) has finite volume.

For an integral ideal, choose an integer basis and form the matrix
of all field embeddings.  The square of its determinant is the
discriminant of that basis; its absolute value is
\(\sqrt{|\operatorname{disc}E|}\,N_E(I)\), by the integral index
formula.  Replacing the two rows \((\sigma,\bar\sigma)\) by
\((\operatorname{Re}\sigma,\operatorname{Im}\sigma)\) divides
the absolute determinant by two.  There are exactly \(d\) pairs,
which proves (RCV.4).  Clearing an integer denominator proves the
fractional-ideal version.  The same argument with only real rows
proves the assertion for \(F\).

Permuting the places or conjugating any of the chosen representatives
only permutes real coordinates or changes the sign of an imaginary
coordinate.  Its absolute determinant is one.  These formulas
therefore hold for the actual embedding and place witnesses in
\(X\).  The zero-dimensional case has the one-point parallelotope,
with volume and empty determinant equal to one.
\end{proof}
% END RCV2

% BEGIN RCV3
\paragraph{RCV3. The actual orthogonal-complement lattice is self-dual.}
\label{claim:RCV3}
Put \(W=\langle e_b,f_b:b<r\rangle^\perp\).  Then
\(\dim_E W=q\), \(L_b=L_0\cap W\), and \(L_b\) is a full
\(\mathcal O_E\)-lattice in \(W\) which is self-dual for \(h|_W\).
\begin{proof}
Pairing a relation among the \(e_b,f_b\) with each \(f_a,e_a\)
proves their linear independence and identifies their Gram matrix
as the hyperbolic matrix.  Their span is nondegenerate.  Thus
\(V=\langle e,f\rangle\perp W\) and
\(\dim_E W=(2r+k+1)-2r=q\).  The projection onto \(W\) is
\[
 p_W(y)=y-\sum_b h(f_b,y)e_b-\sum_b h(e_b,y)f_b.
\tag{RCV.5}
\]
Apply it to the supplied sum decomposition of \(L_0\).  It vanishes
on the two hyperbolic summands and is the identity on \(L_b\),
so its image is precisely \(L_b\).  Pairing with \(e,f\) also
proves directness in (RCV.2), and proves \(L_0\cap W=L_b\).
Finite generation of \(L_0\) and its full span now imply the
same two properties for \(L_b\) in \(W\).

The original \(V\) record supplies Hermitian unimodularity of
\(L_0\).  Its hyperbolic summands are themselves perfect duals:
\(\overline{\mathfrak a_b}=\mathfrak a_b\), and the dual of
\(\mathfrak a_b e_b\) in the \(f_b\) direction is
\(\mathfrak a_b^{-1}f_b\).  More directly, if \(x\in W\)
pairs integrally with \(L_b\), it pairs as zero with all the
hyperbolic summands.  It therefore pairs integrally with \(L_0\),
belongs to \(L_0\) by unimodularity, and belongs to \(L_b\) by
the intersection identity.  The converse follows from integrality
of \(h\) on \(L_0\).  This proves the asserted self-duality,
without choosing a principal \(\mathcal O_E\)-basis for \(L_b\).
\end{proof}
% END RCV3

% BEGIN RCV4
\paragraph{RCV4. The complete trace Gram determinant, in any integer basis.}
\label{claim:RCV4}
For every nonzero fractional ideal \(a\) of \(E\) and every
\(\mathbb Z\)-basis \(m_1,\ldots,m_{2dq}\) of \(aL_b\),
\[
 \left|\det\big(\operatorname{Tr}_{E/\mathbb Q}
                         h(m_i,m_j)\big)\right|
       =|\operatorname{disc}E|^q N_E(a)^{2q}.
\tag{RCV.6}
\]
In particular this proves the full original \texttt{TrDetUnimod}
when \(a=\mathfrak a_b\), for every supplied basis and its original
dimension witness.
\begin{proof}
First let \(L=L_b\), and let \(L^{\mathrm{tr}}\) be its dual for
the rational bilinear form \((x,y)\mapsto\operatorname{Tr}h(x,y)\).
If the trace pairing with every \(y\in L\) is integral, the same
is true with every \(cy\), \(c\in\mathcal O_E\).  By the
definition of the codifferent this says
\(h(x,L)\subset\mathfrak D_{E/\mathbb Q}^{-1}\).  Localize
at a prime of \(\mathcal O_E\).  There the different is principal
and \(L\) is free; RCV3 identifies the Hermitian dual of \(L\)
with \(L\).  Multiplying by a local generator of the different
therefore proves
\[
           L^{\mathrm{tr}}=\mathfrak D_{E/\mathbb Q}^{-1}L.
\tag{RCV.7}
\]
Equality at all primes proves equality of the original fractional
lattices.  Conjugation does not change this formula: the different
is stable under every field automorphism, and we are using an
absolute determinant.

A rank-\(q\) projective \(\mathcal O_E\)-lattice satisfies
\([L:IL]=N_E(I)^q\) for every nonzero integral ideal \(I\).
Indeed, after localization at a prime it is free of rank \(q\),
so the quotient by the \(s\)-th power of that prime has length
\(qs\).  Multiplication over the primes proves the formula.
Apply this to the different in (RCV.7).  Since its norm is
\(|\operatorname{disc}E|\),
\([L^{\mathrm{tr}}:L]=|\operatorname{disc}E|^q\).
In an integer basis of \(L\), the trace-dual basis is obtained
by the inverse Gram matrix.  Thus this index is exactly the
absolute Gram determinant.  This proves (RCV.6) for \(a=\mathcal O_E\).

For integral \(a\), a basis matrix for \(aL\) relative to a
basis of \(L\) has absolute determinant \([L:aL]=N_E(a)^q\).
A bilinear Gram matrix changes by \(B^tGB\), so its determinant
changes by \(N_E(a)^{2q}\).  For arbitrary fractional \(a\),
multiply it by a positive integer to reduce to this case and
cancel the corresponding scalar determinant.  This also proves
that every integer basis has the same answer: a change of such
basis has determinant \(\pm1\).  The rank is \(2dq\) because
the lattice is full in the \(q\)-dimensional \(E\)-space \(W\).
\end{proof}
% END RCV4

% BEGIN RCV5
\paragraph{RCV5. The precise real quadratic form on the row coordinates.}
\label{claim:RCV5}
For a word \(w\), define the additive row map
\[
 \rho_w(x)_v=
 \big(h_v(x,w_{w,v,1}),\ldots,h_v(x,w_{w,v,k}),h_v(x,z_w)\big)
       \quad(x\in W).
\tag{RCV.8}
\]
It extends to a real linear isomorphism
\(W\otimes_{\mathbb Q}\mathbb R\simeq\prod_v\mathbb C^q\).
Its being conjugate-linear in \(x\) over \(E\) is harmless and
is not replaced by complex linearity.  For \(x,y\in W\),
\[
 \frac12\operatorname{Tr}_{E/\mathbb Q}h(x,y)
 =\sum_v\operatorname{Re}\left(
    \sum_{j=1}^k \rho_w(x)_{vj}\overline{\rho_w(y)_{vj}}
       -\frac{\rho_w(x)_{vq}\overline{\rho_w(y)_{vq}}}{T_w}\right).
\tag{RCV.9}
\]
\begin{proof}
The vectors in (RCV.8) have nonsingular Gram matrix
\(G_w=\operatorname{diag}(1_k,-T_w)\) and are \(q\) in number,
so they are a basis of each \(W_v\).  For example the coefficient
expansion is
\[
 y=\sum_j h_v(w_{w,v,j},y)w_{w,v,j}
                         -T_w^{-1}h_v(z_w,y)z_w.
\]
Inserting this into \(h_v(x,y)\) gives the expression inside the
real part in (RCV.9).  The trace of an element of a CM field is
twice the sum of its real parts over one member of each conjugate
pair.  This proves the displayed identity.

The usual real scalar extension of the field-embedding map is an
isomorphism \(E\otimes_{\mathbb Q}\mathbb R\simeq\mathbb C^d\):
its matrix in a rational basis has nonzero discriminant.  Apply this
to coordinates on \(W\), and then apply the nonsingular Hermitian
coordinate maps just described at each place.  This proves the real
linear isomorphism assertion.  In particular the images of any
integer basis of a full lattice in \(W\) are a real basis.
\end{proof}
% END RCV5

% BEGIN RCV6
\paragraph{RCV6. Full row-block covolume.}
\label{claim:RCV6}
For every \(X,w,b\), the source's row lattice is exactly
\(\rho_w(\mathfrak a_b L_b)\).  It is countable and discrete,
has a measurable strict additive fundamental domain, and every
such domain has volume
\[
 c_{R_b}(w)=T_w^d\,2^{-dq}
                \Delta^{q(1-2(r-b))}.
\tag{RCV.10}
\]
This is the full original \texttt{RowBlockCovol}.
\begin{proof}
The map \(\rho_w\) is additive and commutes with integer scalars.
Consequently its image of \(\mathfrak a_bL_b\) is already an
additive subgroup.  The \texttt{AddSubgroup.closure} in the source
means the subgroup generated by that image; it does not mean a
topological closure.  It therefore gives this same lattice.

Choose any integer basis \((m_i)\) of \(\mathfrak a_bL_b\),
and let \(B\) be its real-coordinate matrix under \(\rho_w\).
In real and imaginary coordinates the form on the right of (RCV.9)
has diagonal matrix \(H\), with each entry of
\((1_k,-T_w^{-1})\) repeated twice at each place.  Hence
\(|\det H|=T_w^{-2d}\).  Equations (RCV.6) and (RCV.9) imply
\[
 |\det B|^2 T_w^{-2d}
 =2^{-2dq}|\operatorname{disc}E|^q
                          N_E(\mathfrak a_b)^{2q}.
\]
Taking the positive square root and using (RCV.3) gives precisely
(RCV.10).  The determinant is nonzero; RCV2 now gives countability,
discreteness, existence, and the formula for every measurable
strict domain.  No special choice of row basis, orthonormal frame,
or domain survives in the answer.
\end{proof}
% END RCV6

% BEGIN RCV7
\paragraph{RCV7. The two complex ideal blocks.}
\label{claim:RCV7}
For \(0\leq b<a<r\), the source's \(X\)- and \(K\)-block
lattices and their covolumes are respectively
\[
\begin{array}{c|c|c}
 \text{block}&\text{lattice in }\mathbb C^d&\text{covolume}\\ \hline
 X_{ba}&\sigma(\mathfrak a_b\mathfrak a_a^{-1})
                 &2^{-d}\Delta^{1-2(a-b)}\\
 K_{ba}&\sigma(\mathfrak a_b\mathfrak a_a)
                 &2^{-d}\Delta^{1-2(r-b)-2(r-a)}.
\end{array}
\tag{RCV.11}
\]
Each assertion includes countability, existence of a measurable
strict additive domain, and equality for every such domain.
\begin{proof}
These images are additive subgroups, so again the generated-subgroup
notation does not enlarge them.  Apply (RCV.4) and
\(\sqrt{|\operatorname{disc}E|}=\Delta\).  Multiplicativity of
the fractional ideal norm gives, respectively,
\[
 N_E(\mathfrak a_b\mathfrak a_a^{-1})=\Delta^{-2(a-b)},
 \qquad
 N_E(\mathfrak a_b\mathfrak a_a)
       =\Delta^{-2(r-b)-2(r-a)}.
\]
The ideals are nonzero.  RCV2 supplies all the domain and
countability assertions, with the source's actual embeddings.
\end{proof}
% END RCV7

% BEGIN RCV8
\paragraph{RCV8. The imaginary diagonal block, including the factor at two.}
\label{claim:RCV8}
For \(0\leq a<r\), put
\(I=\mathfrak a_a^2\) and
\(I^- =\{t\in I:\bar t=-t\}\).  The lattice
\(\operatorname{Im}\sigma(I^-)\subset\mathbb R^d\) has covolume
\[
                  c_{D_a}=\Delta^{1/2-2(r-a)}.
\tag{RCV.12}
\]
It is the source's diagonal block, with the same countability,
existence, and every-domain conclusions as in RCV7.
\begin{proof}
Let \(J=\mathfrak D_{F/\mathbb Q}^{-2(r-a)}\), so
\(I=J\mathcal O_E\).  Extension and contraction give
\(I\cap F=J\).  For completeness, this can be checked after
localizing at a prime of \(\mathcal O_F\): write the fractional
ideal as one generator times the local ring, divide by that
generator, and use \(\mathcal O_E\cap F=\mathcal O_F\).

The relative trace carries \(I\) into \(J\).  It is onto,
because \(\theta j\in I\) has trace \(j\) for every \(j\in J\).
We thus have the split exact sequence of integer lattices
\[
 0\longrightarrow I^-\longrightarrow I
       \xrightarrow{\operatorname{Tr}_{E/F}}J\longrightarrow0,
 \qquad j\longmapsto\theta j.
\tag{RCV.13}
\]
This is an integral splitting, including at two.  In the chosen
complex embeddings, every element of \(I^-\) has real part zero,
whereas \(\operatorname{Re}\sigma(\theta j)=\sigma(j)/2\).
Choose integer bases of the kernel and of \(J\), and lift the
latter by \(j\mapsto\theta j\).  The resulting real-coordinate
basis matrix for \(I\) is block triangular.  Its determinant gives
\[
 \operatorname{covol}_{\mathbb C^d}\sigma(I)
 =\operatorname{covol}_{\mathbb R^d}\big(\tfrac12\sigma(J)\big)
       \operatorname{covol}_{\mathbb R^d}\operatorname{Im}\sigma(I^-).
\tag{RCV.14}
\]
Both determinant factors are nonzero; the sequence also shows that
the kernel has integer rank \(d\).  RCV2 and RCV1 give the first
factor on the left as \(2^{-d}\Delta N_F(J)^2\), and the
first factor on the right as \(2^{-d}\sqrt\Delta N_F(J)\).
Their quotient is \(\sqrt\Delta N_F(J)\), which is (RCV.12).
Thus the factors \(2^{-d}\) cancel here.  Replacing this kernel
by the imaginary parts of all of \(I\) would lose exactly the
integral condition that justifies this cancellation.
Finally apply RCV2 to its integer basis to obtain every measurable
strict-domain conclusion.
\end{proof}
% END RCV8

% BEGIN RCV9
\paragraph{RCV9. The complete original block-covolume package.}
\label{claim:RCV9}
For every supplied \(X\), every word \(w\), and every block in
\[
 \mathcal P=(\{(b,a):b<a<r\}\sqcup\{b<r\})
                \sqcup(\{(b,a):b<a<r\}\sqcup\{a<r\}),
\]
the source's \(cLat(X,w,p)\) is a full lattice in its specified
real or complex block space.  Its Lebesgue covolume is exactly
\(cVal(X,w,p)\), as displayed in (RCV.10)--(RCV.12).
This proves \texttt{BlockCovol} with all its original quantifiers.
\begin{proof}
The four alternatives for \(p\) are exactly the \(X\), row,
\(K\), and diagonal blocks proved in RCV6--RCV8.  Their measures
are the standard Lebesgue measures on \(\mathbb C^d\),
\((\mathbb C^q)^d\), \(\mathbb C^d\), and \(\mathbb R^d\),
respectively.  Their orderings differ from real-imaginary coordinate
order only by coordinate permutations, whose absolute determinant
is one.  In each case we have proved all three fields of the original
package: countability of the subgroup, existence of a measurable
strict additive domain, and its stated volume for every supplied
measurable strict additive domain.  Every displayed covolume is
finite and strictly positive.  There is no remaining row-determinant
or diagonal-lattice premise.
\end{proof}
% END RCV9

\subsection{The actual group and its arithmetic subgroup in coordinates}

% BEGIN RCV10
\paragraph{RCV10. A complete polynomial chart for the unipotent group.}
\label{claim:RCV10}
At one real place fix a word \(w\), put
\(G_w=\operatorname{diag}(1_k,-T_w)\), and use the ordered basis
\((e,W_w,f)\) of \(V_v\).  Every element of the original
\(N_v\) has a unique expression
\[
 n(Q,B,K)=
 \begin{pmatrix}
 Q&B&QM\\
 0&1&-G_w^{-1}B^*Q^{-*}\\
 0&0&Q^{-*}
 \end{pmatrix},\qquad
 M=K-\tfrac12Q^{-1}BG_w^{-1}B^*Q^{-*},
\tag{RCV.15}
\]
where \(Q\) is upper unitriangular, \(B\) is an arbitrary
\(r\)-by-\(q\) complex matrix, and \(K^*=-K\).  Here
\(Q^{-*}=(Q^{-1})^*\).
The free coordinates are
\[
 Q_{ba}\ (b<a),\qquad B_{bj}\ (j\leq q),\qquad
 K_{ba}\ (b<a),\qquad \operatorname{Im}K_{aa}.
\tag{RCV.16}
\]
They are exactly the grouped source coordinates \(\operatorname{blk}\);
their real and imaginary parts, with the specified permutation, are
\(\kappa_w\).  They give a homeomorphism from \(N_\infty\) to a
real coordinate space and a measurable section into the ambient group.
Thus the original \texttt{KapSurj} and \texttt{CoordMeas} hold.
\begin{proof}
The source defines \(N_v\) by the isometry condition,
\((n-1)e_a\in\langle e_b:b<a\rangle\), and
\((n-1)e^\perp\subset\langle e\rangle\).  Its additional
condition on \(z_w\) follows from the last one because
\(z_w\perp e\).  In particular the subgroup is independent of
the word.  The first two conditions force the upper block triangular
shape, with upper unitriangular \(Q\) and identity middle block.
Preserving the pairings of the \(e\)- and \(f\)-columns forces
the lower right block to be \(Q^{-*}\).  Pairing the middle
columns with the \(f\)-columns forces the middle right block in
(RCV.15).  Pairing two \(f\)-columns then gives
\[
 M+M^*=-Q^{-1}BG_w^{-1}B^*Q^{-*}.
\tag{RCV.17}
\]
Its solutions are exactly the stated \(M\), with arbitrary
skew-Hermitian \(K\).  Conversely these computations verify every
pairing of the displayed columns, so (RCV.15) is an isometry with
the required flag conditions.  Its diagonal blocks are invertible,
which also verifies that it lies in the actual unitary group.

The coordinate readings are literal: \(Q_{ba}=h(f_b,ne_a)\),
\(B_{bj}=h(f_b,nw_j)\) and \(B_{bq}=h(f_b,nz_w)\), while
\(M=Q^{-1}(h(f_b,nf_a))_{ba}\).  Its skew part is \(K\);
the diagonal entries of the Hermitian correction in (RCV.17)
are real, so \(\operatorname{Im}M_{aa}=\operatorname{Im}K_{aa}\).
These are precisely the four source blocks.

The inverse of an upper unitriangular matrix is the finite polynomial
\(1-U+U^2-\cdots+(-U)^{r-1}\), where \(U=Q-1\).
Consequently (RCV.15) and its coordinate readings on \(N_v\)
are continuous polynomial maps.  The subgroup is closed, since its
flag and isometry equations are closed finite-dimensional equations.
Taking the product over the supplied places yields the stated
homeomorphism and section.  Images of measurable subsets of the
group under this chart are measurable by taking inverse images under
the continuous inverse chart.  On ambient matrices, the source uses
the total matrix-inverse operation in the reading of \(M\); it is
Borel measurable, being the adjugate times the reciprocal determinant
with the specified value at determinant zero.  This verifies the
ambient measurability required by \texttt{CoordMeas} as well.
\end{proof}
% END RCV10

% BEGIN RCV11
\paragraph{RCV11. The full triangular multiplication law and its measure.}
\label{claim:RCV11}
Give the four blocks the source heights
\[
 h(X_{ba})=a-b,\quad h(R_b)=2r-b,\quad
 h(K_{ba})=4r+1-b-a,\quad h(D_a)=4r+1-2a.
\tag{RCV.18}
\]
In block coordinates multiplication has the form
\[
       (x*y)_p=x_p+y_p+\mathcal B_p(x_{<h(p)},y_{<h(p)}),
\tag{RCV.19}
\]
where each correction is a polynomial, depends only on strictly
lower-height blocks in both arguments, and vanishes if either
argument is zero.  This is the full original \texttt{KapLaw}.
Coordinate Lebesgue measure is both left and right invariant.
\begin{proof}
Multiply (RCV.15) for \((Q,B,K)\) and \((Q',B',K')\).
Writing \(U=Q^{-1}\), \(U'=Q'^{-1}\), one obtains
\[
 \begin{split}
 Q''&=QQ',\qquad B''=QB'+B,\\
 M''&=M'+U'MU'^*-U'UBG_w^{-1}B'^*U'^*,\\
 K''&=K'+U'KU'^*-\operatorname{skew}
                         (U'UBG_w^{-1}B'^*U'^*),
 \end{split}
\tag{RCV.20}
\]
where \(\operatorname{skew}(A)=(A-A^*)/2\).  The last line
uses that \(U'HU'^*\) is Hermitian whenever \(H\) is Hermitian,
so the Hermitian part of \(M\) contributes no skew part.

In \(Q''_{ba}\), every correction has indices \(b<c<a\),
so its two \(X\)-heights are strictly less than \(a-b\).
In row \(b\) of \(B''\), every correction uses \(Q_{bc}\)
and row \(c>b\) of \(B'\), again of strictly lower heights.
For the central block, every entry of \(U'KU'^*\) contributing
to position \((b,a)\) has central indices \(c\geq b,e\geq a\).
Apart from \((c,e)=(b,a)\), which contributes exactly \(K_{ba}\),
its central height is \(4r+1-c-e<4r+1-b-a\).  This statement
also covers conjugate lower entries and diagonal entries, since their
heights use the sum of the indices.  Entries of \(U,U'\) are
polynomials in the lower-height \(X\)-blocks.  Every row block
appearing in the remaining term in (RCV.20) has height at most
\(2r\), whereas every central height is at least \(2r+3\).
This proves strict dependence for both off-diagonal and diagonal
central coordinates.  All the corrections are polynomials, hence
jointly measurable.  Substitution of the identity element proves
their asserted vanishing; strict dependence makes the same statement
true whenever just the relevant lower coordinates vanish.

Order the finitely many blocks by increasing height, breaking ties
arbitrarily.  Left translation by a fixed point is a triangular
map with identity diagonal blocks: successively it translates each
block by a function of the preceding blocks.  Each such map is a
global measurable bijection and preserves product Lebesgue measure,
by the translation invariance in that block and Tonelli's theorem.
Their finite composition has the same property.  The identical
argument with the other argument fixed proves right invariance.
Thus the coordinate measure is a genuine Haar measure on the actual
closed unipotent subgroup.
\end{proof}
% END RCV11

% BEGIN RCV12
\paragraph{RCV12. The needed finite-level membership, on the supplied lattice.}
\label{claim:RCV12}
For every supplied \(X,w\) and every finite adelic element
\(n\in N_{\mathrm{fin}}(w)\cap K\),
\[
                         x_w^{-1}nx_w\in K_2.
\tag{RCV.21}
\]
Here \(K\), \(K_2\), \(x_w=X.xu(w)\), and all finite frame
vectors are the original ones.  This is the full source theorem
\texttt{unipotent\_conj\_mem\_level}; it assumes no character
identity.
\begin{proof}
Let \(A\) be the integral finite adelic coefficient ring and
\(L=X.B.I.Lf\).  Its form is integral.  Put \(g=X.B.g\).
The supplied \texttt{basisApprox} and \texttt{xu\_z} assertions
give one positive integer \(C\), coprime to \(g\), such that
\[
 \begin{split}
 C(e_b-e_b^0),\ C(f_b-f_b^0)&\in gL,\\
 C(x_w^{-1}z_w-z^0)&\in gL
 \end{split}
\tag{RCV.22}
\]
for all \(b\).  Take the product of their finitely many individual
clearing integers.  Work for the moment in \(L_C=L[1/C]\) over
\(A_C=A[1/C]\), modulo \(gL_C\).  The actual \(e_b,f_b\)
belong to \(L_C\) and are congruent to \(e_b^0,f_b^0\).
Moreover, the natural map \(L/gL\to L_C/gL_C\) is an
isomorphism: multiplication by \(C\) is invertible modulo \(g\)
by an integer Bezout identity.  This observation will descend every
congruence below to the original module, with no localization left
in the conclusion.

Set \(a=x_w^{-1}nx_w\).  The supplied \(x_w\in K\) and
\(n\in K\) imply \(a\in K\); all three maps preserve \(L_C\).
The supplied \(x_w\) fixes each \(e_b,f_b\), hence preserves
their orthogonal complement.  Conjugation by it therefore preserves
the flag unipotent conditions.  In particular
\((a-1)e_b\in\langle e_c:c<b\rangle\).
An element \(s\in L_C\cap\langle e_c:c\in S\rangle\)
has the exact expansion
\(s=\sum_{c\in S}h(f_c,s)e_c\), whose coefficients lie in
\(A_C\).  Modulo \(g\), its image belongs to the span of
\(e_c^0\), \(c\in S\).  Apply this with
\(s=(a-1)e_b\) and then use \(e_b\equiv e_b^0\) to obtain
the first family of congruences defining \(K_2\).

Next, \((n-1)z_w\in\langle e\rangle\).  Since \(x_w\)
fixes that span, \((a-1)x_w^{-1}z_w\) belongs to it and to
\(L_C\).  Equation (RCV.22) gives
\((a-1)z^0\in\langle e^0\rangle+gL_C\), the second
congruence of \(K_2\).

Finally let \(\beta_i\) be any supplied \(bw_i\).  Its original
Gram congruences give \(h(e_b^0,\beta_i)\in gA\), and hence
\(h(e_b,\beta_i)\in gA_C\).  Decompose
\[
 \beta_i=\sum_b h(f_b,\beta_i)e_b+p_W\beta_i
                              +\sum_b h(e_b,\beta_i)f_b.
\]
All these terms belong to \(L_C\), because the frame vectors
and their pairings are integral there.  The last sum belongs to
\(gL_C\).  The flag conditions put \((a-1)\) of the first
two terms in \(\langle e\rangle\); the first term uses the
upper flag condition and the second uses \(p_W\beta_i\in e^\perp\).
Thus \((a-1)\beta_i\) is, modulo \(gL_C\), an integral
combination of \(e_b^0\).  This is stronger than the third
congruence of \(K_2\), which permits earlier \(bw_j\) as well.

All the vectors \((a-1)e_b^0,(a-1)z^0,(a-1)\beta_i\)
belong to \(L\).  The isomorphism modulo \(g\) descends the
three span congruences to \(L/gL\): coefficients in \(A_C/gA_C\)
come from \(A/gA\).  Together with \(a\in K\), they are
exactly \(a\in K_2\).  Notice that the argument uses only
membership, so neither \(\mathrm{PsiDd}\) nor any equality of
finite characters has been introduced as an extra premise.
\end{proof}
% END RCV12

% BEGIN RCV13
\paragraph{RCV13. The exact rational-arithmetic dictionary.}
\label{claim:RCV13}
For every \(X,w\), the subgroup \(\Gamma_w=\mathrm{NOinf}(X,w)\)
is the archimedean image of
\[
            \{n\in N(E):nL_0=L_0\}.
\tag{RCV.23}
\]
The rational isometry and its finite and infinite components are
linked by the original \texttt{RatLink}.  This proves the complete
original \texttt{NOratC}, not only an inclusion in an auxiliary
arithmetic group.
\begin{proof}
By definition, \(s\in\Gamma_w\) means there is an original
rational adelic point \(\gamma\in X.B.I.GFset\) with
\(\gamma_\infty=s\), whose two components lie in the specified
unipotent groups, and such that
\(x_w^{-1}\gamma_f x_w\in K_2\).  Since \(K_2\subset K\)
and \(x_w\in K\), we have \(\gamma_f\in K\).

For a rational vector \(y\), membership in \(L_0\) is equivalent
to membership of its finite adelic image in \(L\).  One elementary
verification chooses an integer basis of the full lattice \(L_0\).
The rational coordinates of \(y\) are integral in every \(\mathbb Z_p\)
if and only if they are integers.  Scalar extension to the finite
adeles is exactly the original \(L\), by the supplied lattice
and coefficient-coordinate identities.  Consequently preservation of
\(L\) by a rational isometry and its inverse is equivalent to
preservation of \(L_0\).

The original rational point therefore gives an isometry \(n\)
preserving \(L_0\).  Its infinite flag conditions descend to
rational ones.  Indeed, decompose a rational vector using the
hyperbolic projection (RCV.5); vanishing of any unwanted coefficient
at an embedding forces its vanishing in \(E\), and the same
argument applied to a rational basis of \(e^\perp\) proves
the universal perpendicular-space condition.  Thus \(n\in N(E)\).
This proves one direction of (RCV.23).

Conversely, a rational isometry in (RCV.23) gives its actual pair of
adelic components through the supplied coordinates.  The isometry,
flag, and perpendicular-space equations extend to all completions;
the latter can be checked on a rational basis of \(e^\perp\).
Hence its components belong to \(N_\infty\) and \(N_{\mathrm{fin}}\).
The lattice argument above gives \(\gamma_f\in K\), and RCV12
gives \(x_w^{-1}\gamma_f x_w\in K_2\).  These are precisely
the defining witnesses for \(\Gamma_w\).  The coefficient map
is injective, so no rational point is lost or duplicated.
In particular \(\Gamma_w\) is countable and is a subgroup of
the actual \(N_\infty\).  The source's extra condition on
\(z_w\) in \(N(E)\) is redundant for the same reason as in
RCV10, so this dictionary also respects every word witness.
\end{proof}
% END RCV13

% BEGIN RCV14
\paragraph{RCV14. Rational coordinates and reconstruction.}
\label{claim:RCV14}
For \(n\in N(E)\), define
\[
 Q_{ba}=h(f_b,ne_a),\quad A_{ba}=h(f_b,nf_a),\quad
 M=Q^{-1}A,\quad x_b=p_W(n^{-1}f_b).
\tag{RCV.24}
\]
Then \(Q\) is upper unitriangular, \(x_b\in W\), and, putting
\(G(x)_{ba}=h(x_b,x_a)\),
\[
               M+M^*=-H,\qquad H=Q^{-1}G(x)Q^{-*}.
\tag{RCV.25}
\]
Conversely every such rational tuple \((Q,x,M)\) reconstructs a
unique element of \(N(E)\).  At every chosen place its block
coordinates are \(\sigma(Q_{ba})\), \(\rho_w(x_b)\), the
upper entries of \(\operatorname{skew}\sigma(M)\), and
\(\operatorname{Im}\sigma(M_{aa})\).
This proves the original \texttt{DictGram}, \texttt{DictEx}, and
\texttt{DictPlace}.
\begin{proof}
The group conditions give the assertions about \(Q\).  For
\(y\in W\), preservation of the form gives
\[
 h(f_b,ny)=h(n^{-1}f_b,y)=h(x_b,y).
\]
Thus \(ny=y+\sum_bh(x_b,y)e_b\).  Put \(U=Q^{-*}\).
The complete action on the three rational summands is forced to be
\[
 \begin{split}
 ne_a&=\sum_b Q_{ba}e_b,\\
 ny&=y+\sum_b h(x_b,y)e_b\quad(y\in W),\\
 nf_a&=\sum_b(QM)_{ba}e_b
                         -\sum_b U_{ba}x_b+\sum_b U_{ba}f_b.
 \end{split}
\tag{RCV.26}
\]
The middle component of the last line is obtained by pairing with
all vectors \(ny\), \(y\in W\); nondegeneracy on \(W\)
makes it unique.  Pairing two last-line columns gives
\(M+M^*+U^*G(x)U=0\), which is (RCV.25).

Conversely, (RCV.26) defines a rational linear map for every tuple
as in the statement.  Its three diagonal blocks are \(Q,1,U\),
so it is invertible.  Pairings of its \(e\)-columns and its
\(W\)-columns are unchanged; their pairings with the \(f\)-columns
are unchanged by the two correction terms in the last line.
The remaining \(f,f\) pairing is zero precisely by (RCV.25).
Thus it is an isometry and has the required unipotent flag behavior.
For this map, pairing \(n^{-1}f_b\) with any \(y\in W\)
recovers \(h(x_b,y)\), so its projected vector is indeed \(x_b\).
All its readings (RCV.24) are therefore the prescribed tuple, proving
uniqueness and the exact reconstruction assertion.

Finally extend (RCV.26) to each chosen complex embedding.  Its row
entries on the frame of \(W_v\) are \(\rho_w(x_b)\).
The remaining readings agree directly with (RCV.15)--(RCV.17).
This proves the exact placement statement on the original rational
and adelic components, including the order of conjugation in the
row map.
\end{proof}
% END RCV14

% BEGIN RCV15
\paragraph{RCV15. Exact preservation of the original lattice.}
\label{claim:RCV15}
For a tuple in RCV14, the associated \(n\) preserves \(L_0\)
if and only if
\[
 \begin{split}
 Q_{ba}&\in\mathfrak a_b\mathfrak a_a^{-1}\quad(b<a),\\
 x_b&\in\mathfrak a_b L_b,\\
 M_{ba}&\in\mathfrak a_b\mathfrak a_a\quad(a,b<r).
 \end{split}
\tag{RCV.27}
\]
This is the full original \texttt{DictLat}.
\begin{proof}
Suppose first that \(nL_0=L_0\).  Apply \(n\) to
\(c e_a\), \(c\in\mathfrak a_a\), and extract the \(e_b\)
coefficient by pairing with \(f_b\).  Equation (RCV.2) gives
\(Q_{ba}\mathfrak a_a\subset\mathfrak a_b\), equivalently
the first line of (RCV.27).  Apply \(n^{-1}\) to
\(c f_b\), \(c\in\mathfrak a_b^{-1}\), and project onto
\(W\).  This gives \(\mathfrak a_b^{-1}x_b\subset L_b\),
equivalently \(x_b\in\mathfrak a_bL_b\); the equivalence may
be checked after localizing the invertible ideal at each prime.
Likewise applying \(n\) to \(\mathfrak a_a^{-1}f_a\)
gives \(A_{ba}\in\mathfrak a_b\mathfrak a_a\).
The finite series for \(Q^{-1}\) has entries
\((Q^{-1})_{bc}\in\mathfrak a_b\mathfrak a_c^{-1}\)
for \(b\leq c\), because the ideals telescope along each
matrix-product path.  Hence \(M=Q^{-1}A\) satisfies the last
line of (RCV.27).

Conversely assume (RCV.27).  The same inverse-series argument
gives the stated ideal bounds on \(Q^{-1}\), and therefore
\(U_{ba}\in\mathfrak a_b^{-1}\mathfrak a_a\) whenever it
is nonzero.  Formula (RCV.26) now preserves each lattice generator
in (RCV.2).  For \(y\in L_b\), self-duality and conjugation
invariance give \(h(x_b,y)\in\mathfrak a_b\).
For \(c f_a\), \(c\in\mathfrak a_a^{-1}\), its three
components have respectively coefficients in \(\mathfrak a_b\),
vectors in \(L_b\), and coefficients in \(\mathfrak a_b^{-1}\):
use \((QM)_{ba}\in\mathfrak a_b\mathfrak a_a\) and
\(U_{ba}x_b\in\mathfrak a_aL_b\).  Thus \(nL_0\subset L_0\).

Because \(n\) is an isometry and \(L_0\) is self-dual,
\((nL_0)^{\#}=nL_0\), where \(\#\) denotes Hermitian
duality.  Taking duals reverses the inclusion and gives
\(L_0\subset nL_0\).  Hence it is equality.  This final step
proves preservation, rather than only integral inclusion, for the
actual possibly nonprincipal lattice.
\end{proof}
% END RCV15

% BEGIN RCV16
\paragraph{RCV16. The complete shifted-fiber description of the arithmetic subgroup.}
\label{claim:RCV16}
In the block coordinates of RCV10, \(\Gamma_w\) has exactly the
source's \texttt{NOfib} description
\[
 t\in\operatorname{blk}_w(\Gamma_w)
 \quad\Longleftrightarrow\quad
 t_p-\tau_p(t_{<h(p)})\in cLat(X,w,p)\quad\hbox{for every }p.
\tag{RCV.28}
\]
The shifts of the \(X\) and row blocks are zero.  On every real
coordinate tuple form the upper unitriangular \(Q\) from its
\(X\)-blocks and the matrix \(B\) from its row blocks, and put,
at each place, \(H=Q^{-1}BG_w^{-1}B^*Q^{-*}\).
When these blocks are rationally recoverable, this is exactly
\(\sigma(Q^{-1}G(x)Q^{-*})\) by RCV5.  The other shifts are
\[
 \tau_{K_{ba}}=\tfrac12 H_{ba},\qquad
 \tau_{D_a}=-\operatorname{Re}H_{aa}
                                  \operatorname{Im}\sigma(\theta).
\tag{RCV.29}
\]
They are the literal \(\mathrm{tauQ}\) of the source, depend only
on lower-height blocks, and vanish at zero.  Both directions of
(RCV.28), including existence of a rational lift of every admissible
coordinate tuple, hold.  More precisely, for every rational
\(\mathrm{DataOK}\) tuple the shifted-lattice conditions on its
\(\mathrm{blkQ}\) are equivalent to \(\mathrm{LatQ}\), and
every real tuple satisfying those conditions is the \(\mathrm{blkQ}\)
of a rational \(\mathrm{DataOK}\) tuple.  These are the full
original \texttt{FibQ} and \texttt{FibOnto}.
\begin{proof}
The ideal bounds of (RCV.27) imply
\(G(x)_{ba}\in\mathfrak a_b\mathfrak a_a\), and the
inverse-matrix bounds used in RCV15 then imply the same bound for
\(H_{ba}\).  From (RCV.25),
\[
 \operatorname{skew}(M)_{ba}=M_{ba}+H_{ba}/2.
\tag{RCV.30}
\]
Thus the upper off-diagonal lattice condition on \(M\) is
exactly the \(K\)-condition in (RCV.28).  Its lower entries are
forced by \(M_{ab}=-H_{ab}-\overline{M_{ba}}\), and obey
the corresponding ideal condition as well.

On the diagonal, \(H_{aa}\in F\), and
\[
                  \xi_a=M_{aa}+H_{aa}\theta
\]
is anti-invariant because \(M_{aa}+\bar M_{aa}=-H_{aa}\)
and \(\theta+\bar\theta=1\).  It lies in \(\mathfrak a_a^2\)
exactly when \(M_{aa}\) does, since \(\theta\) is integral.
Its imaginary coordinate is
\(\operatorname{Im}\sigma(M_{aa})+
 \operatorname{Re}\sigma(H_{aa})\operatorname{Im}\sigma(\theta)\).
This gives exactly the diagonal shift (RCV.29), with its stated
sign.  RCV13--RCV15 prove the forward implication.

For the converse, the unshifted \(X\) blocks uniquely recover
rational entries of \(Q\) in the specified ideals.  The row
blocks uniquely recover \(x_b\in\mathfrak a_bL_b\), by
injectivity of (RCV.8).  They give the same rational \(H\) as
above.  For each \(b<a\), the shifted \(K\)-coordinate
recovers a unique \(M_{ba}\in\mathfrak a_b\mathfrak a_a\).
Set \(M_{ab}=-H_{ab}-\overline{M_{ba}}\).  Each diagonal
condition recovers an anti-invariant \(\xi_a\in\mathfrak a_a^2\);
set \(M_{aa}=\xi_a-H_{aa}\theta\).  These formulas give
every entry of a rational matrix satisfying (RCV.25) and (RCV.27).
RCV14 reconstructs its isometry, RCV15 proves lattice preservation,
and RCV13 places its actual infinite component in \(\Gamma_w\).
Its coordinates are exactly the supplied tuple, by (RCV.30) and
the diagonal computation.  This proves surjectivity onto the entire
shifted-fiber set, not merely onto coordinates of points already
known to be rational.

For the stated \texttt{FibQ} equivalence, a rational tuple failing
an \(X\)- or row-lattice condition fails the same unshifted
condition on both sides.  Otherwise the computations (RCV.30) and
the diagonal identity prove the equivalence entry by entry, for
every \(\mathrm{DataOK}\) tuple.  The reconstruction in the
preceding paragraph is exactly the universal \texttt{FibOnto}
return.  All \(X\)- and row heights are below all central heights.
Hence the actual polynomial shifts have the required strict
lower-height dependence.  They are measurable as well, although
the original \texttt{NOfib} does not require this extra property.
Finally, subgroup closure and inclusion in \(N_\infty\) were
proved in RCV13.  These supply the other fields of \texttt{NOfib}.
\end{proof}
% END RCV16

\subsection{All fundamental domains and all supplied Haar measures}

% BEGIN RCV17
\paragraph{RCV17. A triangular lattice has the product covolume.}
\label{claim:RCV17}
For the actual group of RCV11 and subgroup of RCV16, choose any
half-open fundamental parallelotope \(P_p\) for each block lattice.
The Cartesian product \(P=\prod_p P_p\) is a strict fundamental
domain in block coordinates.  Every measurable strict domain
\(A\) for this subgroup has
\[
                     \operatorname{vol}A=\prod_p cVal(X,w,p).
\tag{RCV.31}
\]
Consequently, for every original ambient
\(\operatorname{IsSFund}(\Gamma_w,\Omega)\),
\(\kappa_w(\Omega\cap N_\infty)\) is measurable and satisfies
(RCV.31).
\begin{proof}
Order the blocks by increasing height.  The shifted-fiber condition
and its surjectivity in RCV16 show that, for every block \(p\)
and every \(\lambda\in cLat_p\), there is \(\gamma\in\Gamma_w\)
whose preceding blocks are zero and whose \(p\)-th block is
\(\lambda\).  To construct it, set those earlier coordinates to
zero, the current coordinate to \(\lambda\), and choose each
later coordinate in its nonempty shifted lattice.  All earlier
conditions hold because the shifts vanish at zero.

Given an arbitrary point of the real coordinate space, reduce its
first coordinate into \(P_1\) by left multiplication by such a
\(\gamma\).  At the next step use an element with all preceding
coordinates zero.  Equation (RCV.19) and its zero-argument property
say that this leaves those coordinates fixed and translates the
current one by exactly \(\lambda\).  After finitely many steps
the point lies in \(P\).  For uniqueness, suppose \(x,y\in P\)
and \(y=\gamma*x\).  At the first coordinate the difference
is in its lattice, so the strict parallelotope condition gives
\(\gamma_1=0\) and \(x_1=y_1\).  Inductively the earlier
coordinates of \(\gamma\) vanish, so its current coordinate
belongs to the unshifted lattice and the correction in (RCV.19)
vanishes.  The same argument gives its vanishing.  Thus
\(\gamma=1\) and \(x=y\).  This proves that \(P\) is a
strict domain even though the subgroup itself has shifted fibers.

Its product Lebesgue volume is \(\prod_pcVal_p\), by RCV9 and
Tonelli.  The subgroup is countable by RCV13, and left translations
preserve Lebesgue measure by RCV11.  Partition any other measurable
strict domain \(A\) by its intersections with the translates of
\(P\), and translate the pieces back.  They partition \(P\),
exactly as in RCV2.  This proves (RCV.31) before assuming finiteness
of \(A\).

An original \texttt{IsSFund} is a measurable set in the ambient
archimedean group; its defining translates cover that whole group
and meet it uniquely.  Because \(\Gamma_w\subset N_\infty\),
its intersection with \(N_\infty\) has the same strict-domain
properties on \(N_\infty\).  Its image is measurable by the
inverse chart in RCV10.  Finally \(\kappa_w\) only separates
complex entries into real and imaginary coordinates and permutes
them.  It preserves precisely the block product Lebesgue measure.
Thus (RCV.31) applies to every original ambient domain witness.
\end{proof}
% END RCV17

% BEGIN RCV18
\paragraph{RCV18. Multiplication of the four exact block constants.}
\label{claim:RCV18}
Put
\[
 e_N=-\frac{8r^3+6r^2+6kr^2+r}{6},\qquad
 C_w=T_w^{rd}2^{-d(rk+r^2)}\Delta^{e_N}.
\tag{RCV.32}
\]
Then \(\prod_p cVal(X,w,p)=C_w\), and hence the full
original \texttt{CoordCovol} holds with its literal return
\(\mathrm{covolN}(X,w)/\mathrm{twoU}=C_w\).
\begin{proof}
Only the \(r\) row blocks contain \(T_w\); their product gives
\(T_w^{rd}\).  There are two families of \(r(r-1)/2\)
complex ideal blocks, and \(r\) row blocks of complex dimension
\(dq\).  Thus the power of two in the product is
\(-d(r(r-1)+rq)=-d(rk+r^2)\).

We record the discriminant sum, so that no root count is hidden.
Using \(\sum_{a=0}^{r-1}a=r(r-1)/2\) and
\(\sum_{a=0}^{r-1}a^2=r(r-1)(2r-1)/6\), the four sums are
\[
\begin{array}{c|c}
 X&-r(r-1)(2r-1)/6\\
 R&-qr^2\\
 K&r(r-1)(2r-1)/2-2r^2(r-1)\\
 D&-r^2-r/2.
\end{array}
\tag{RCV.33}
\]
For example the first is
\(\sum_{a=0}^{r-1}\sum_{b<a}(1-2(a-b))\); the third is
\(\sum_{a=0}^{r-1}\sum_{b<a}(1-4r+2a+2b)\).
Adding the four displayed sums and substituting \(q=k+1\)
gives exactly \(e_N\).  They are identities also at \(r=0\).
Positivity of \(T_w,2,\Delta\) permits ordinary addition of
the real exponents.

The source defines
\[
 \mathrm{covolN}(X,w)
 =T_w^{rd}2^{-d(rk+r(r-1)/2)}\Delta^{e_N},\qquad
 \mathrm{twoU}=2^{dr(r+1)/2}.
\]
Since \(r(r-1)/2+r(r+1)/2=r^2\), their quotient is \(C_w\).
Combine this arithmetic with RCV17, including its measurability
and every-domain statements, to obtain the complete original
\texttt{CoordCovol}.  Its right side is finite and strictly positive.
\end{proof}
% END RCV18

% BEGIN RCV19
\paragraph{RCV19. The original Haar formula with every normalization.}
\label{claim:RCV19}
Let \(\mu\) be any ambient measure satisfying the original
\(\mathrm{IsHaarOnSet}(N_\infty,\mu)\): it is sigma-finite,
nonzero, supported on \(N_\infty\), and invariant under its
left translations.  Let \(\mathrm{Box}_w\) be the original
\(\mathrm{boxN}(X,w)\), obtained by placing every real coordinate
of \(\kappa_w\) in \([0,1)\).  Then
\[
 (\kappa_w)_*\mu=\mu(\mathrm{Box}_w)\,d\kappa,
 \qquad 0<\mu(\mathrm{Box}_w)<\infty.
\tag{RCV.34}
\]
For every original \(\operatorname{IsSFund}(\Gamma_w,\Omega)\),
\[
                 \mu(\Omega)=C_w\,\mu(\mathrm{Box}_w).
\tag{RCV.35}
\]
These are the full \texttt{CoordHaar} and \texttt{CovolN},
with their original \([0,\infty]\)-valued equalities understood
via \(\operatorname{ofReal}\) on the finite constants.
\begin{proof}
The coordinate Lebesgue measure transported through RCV10 is a
nonzero locally finite Haar measure by RCV11.  The subgroup is
closed and second countable.  AMD8 applies at exactly the weak
\texttt{IsHaarOnSet} scope above: restriction to the supported
closed subgroup is a genuine Haar measure, and uniqueness gives
one finite strictly positive scalar multiplying the coordinate
measure.  The coordinate unit cube has volume one, including
in dimension zero.  Evaluating at its inverse image identifies
that scalar as \(\mu(\mathrm{Box}_w)\), proving (RCV.34).
Support on \(N_\infty\) justifies this assertion for the ambient
map \(\kappa_w\), including its values off the subgroup.

The same support gives \(\mu(\Omega)=\mu(\Omega\cap N_\infty)\).
Apply (RCV.34) and RCV18 to the measurable coordinate image of
this set.  This gives (RCV.35).  In particular its finiteness and
positivity are conclusions.  The proof has fixed neither a Haar
normalization nor a fundamental domain; it applies to every
witness in the original statement.
\end{proof}
% END RCV19

\subsection{The prescribed Lie measure and the independent reference word}

% BEGIN RCV20
\paragraph{RCV20. The Euclidean Lie-algebra volume factor.}
\label{claim:RCV20}
For arbitrary nonnegative integers \(r,k\), take the source's
standard group \(N_{k,r}\) with middle Gram matrix
\(D_0=\operatorname{diag}(1_k,-1)\).  Its actual realified
Lie space consists exactly of
\[
 L(A,B,K)=
 \begin{pmatrix}
 A&B&K\\ 0&0&-D_0B^*\\0&0&-A^*
 \end{pmatrix},\qquad
 A\text{ strictly upper triangular},\quad K^*=-K.
\tag{RCV.36}
\]
If the independent complex coordinates are the upper entries of
\(A,K\) and all entries of \(B\), and the real coordinates
are \(\operatorname{Im}K_{aa}\), its Euclidean volume is
\[
        2^{r^2+rk}\,L_*\!\left(dA\,dB\,dK_{\mathrm{upper}}
                                      \,d\operatorname{Im}K_{\mathrm{diag}}\right).
\tag{RCV.37}
\]
This proves the complete original \texttt{ConstStep}.
\begin{proof}
Differentiate the standard upper-flag and isometry equations for
\(\exp(tZ)\) at zero.  They force (RCV.36).  Conversely a matrix
of that form is nilpotent and skew-adjoint for the standard
hyperbolic Hermitian matrix.  Its finite exponential preserves
that form, preserves the upper flag and acts trivially on its
specified quotients.  Thus \(\exp(tZ)\in N_{k,r}\) for every
real \(t\).  This proves equality with the source's
\(\mathrm{lieOf}(N_{k,r})\), and hence with its real linear
span \(\mathrm{lieE}(N_{k,r})\); no larger tangent space has
been substituted.

The source's ambient Euclidean norm is the sum of the squares of
the real and imaginary parts of every matrix entry.  Therefore
\[
 \|L(A,B,K)\|^2
  =2\left(\sum_{b<a}|A_{ba}|^2+
           \sum_{b,j}|B_{bj}|^2+
           \sum_{b<a}|K_{ba}|^2\right)
                          +\sum_a(\operatorname{Im}K_{aa})^2.
\tag{RCV.38}
\]
Indeed the \(A\) entries occur again in \(-A^*\), the \(B\)
entries again in \(-D_0B^*\), and every off-diagonal entry of
\(K\) has its conjugate partner.  A diagonal imaginary entry
occurs once.  Polarization shows that distinct real coordinate
directions are orthogonal, with length \(\sqrt2\) for each
real or imaginary part of a complex parameter and length one
for a diagonal parameter.  The number of complex parameters is
\[
              2\binom r2+r(k+1)=r^2+rk.
\]
The real volume factor is consequently
\((\sqrt2)^{2(r^2+rk)}=2^{r^2+rk}\), proving (RCV.37).
The coordinate linear map is a measurable isomorphism onto the
actual Lie space, as the entries in (RCV.36) recover every parameter.
This supplies both the map and the volume equality demanded by
\texttt{ConstStep}, including the empty-coordinate case.
\end{proof}
% END RCV20

% BEGIN RCV21
\paragraph{RCV21. The exponential-to-polynomial chart has determinant one.}
\label{claim:RCV21}
Let \(\Phi\) denote (RCV.15) in the standard normalized frame
\(T=1\).  Then the source's prescribed measure satisfies
\[
 \mathrm{haarLie}(N_{k,r},1)
   =2^{r^2+rk}\,\Phi_*d\kappa_0
                     \quad(r,k\geq0).
\tag{RCV.39}
\]
In particular the full original \texttt{ShearStep} and
\texttt{ExpChart} hold, with no free proportionality constant.
\begin{proof}
Assign degrees \(b\), \(r\), and \(2r-b\) to the standard
basis vectors \(e_b,W_j,f_b\), respectively.  A potentially
nonzero entry of (RCV.36) increases this degree strictly from its
row to its column.  Therefore it is nilpotent, and its exponential
and logarithm are finite polynomial sums.  Assign the independent
coordinates the corresponding positive weights
\[
 a-b\text{ for }A_{ba},\quad r-b\text{ for }B_{bj},\quad
 2r-b-a\text{ for }K_{ba},\quad 2r-2a\text{ on the diagonal}.
\tag{RCV.40}
\]
Every term of degree at least two in an exponential entry follows
a path of at least two positive steps.  Each factor therefore has
strictly smaller weight than the resulting entry.  It follows that
the \(Q\)- and \(B\)-coordinates read from \(\exp L\) are
their linear parameters plus polynomials in smaller-weight parameters.
For the central coordinate, reading \(M=Q^{-1}C\) introduces
only products of positive-weight entries.  Subtracting its Hermitian
part, equivalently reading \(\operatorname{skew}(M)\), again
introduces only such products.  Its linear term is exactly \(K\),
and on the diagonal it is exactly its imaginary parameter.

Thus the map \(s=\Phi^{-1}\circ\exp\circ L\), written in
real and imaginary coordinates and ordered by (RCV.40), has the form
\(s_j(t)=t_j+P_j(t_{\text{smaller weight}})\).  This statement
also holds across equal weights, since every nonlinear factor has
strictly smaller weight.  Successive subtraction gives a polynomial
inverse.  As in RCV11, repeated translation in one coordinate block
and Tonelli show that \(s\) preserves Lebesgue measure exactly.
It is the measurable, measure-preserving map required by
\texttt{ShearStep}; by construction \(\Phi(s(t))=\exp L(t)\).

The definition of \(\mathrm{haarLie}(N_{k,r},1)\) is the
pushforward of the Euclidean volume of its actual realified Lie
space under matrix exponential.  Insert (RCV.37), then use
\(\exp L=\Phi\circ s\) and \(s_*d\kappa_0=d\kappa_0\).
This gives (RCV.39) with the exact coefficient shown.  It also
explains why replacing the exponential chart by the polynomial
chart does not introduce an additional Jacobian factor.
\end{proof}
% END RCV21

% BEGIN RCV22
\paragraph{RCV22. The full local reference-measure formula.}
\label{claim:RCV22}
Take any integers \(n,r,k\) with \(n=2r+k+1\), any signature
\(\epsilon\), and any supplied frame \(e,f,w_1,\ldots,w_k,z\)
in \(\mathbb C^n\) satisfying the hyperbolic pairings,
orthonormal positive \(w_j\), all stated orthogonalities, and
\(h_\epsilon(z,z)=-T\) with \(T>0\).
Let \(\nu\) be the literal source measure
\(\mathrm{refNg}(n,\epsilon,e,f,w,z/\sqrt T)\).
Then the full source coordinate map \(\kappa_{e,f,w,z}\) is
\(\nu\)-almost-everywhere measurable and
\[
       (\kappa_{e,f,w,z})_*\nu
                =2^{r^2+rk}T^{-r}\,d\kappa.
\tag{RCV.41}
\]
Furthermore \(\nu\) is an original \texttt{IsHaarLoc} measure
on \(N(\epsilon,e,z')\) for every \(z'\perp e\), without
a norm or nonzero assumption on \(z'\).
\begin{proof}
The normalized frame \((e,w,z/\sqrt T,f)\) has the source's
standard Gram matrix.  Its column matrix is invertible by
nondegeneracy of that matrix.  The corresponding conjugation
\(g\mapsto P^{-1}gP\) is a closed measurable embedding of the
actual unitary group into standard matrices, with the expected
inverse on its image.  The source defines \(\nu\) as the
pullback of \(\mathrm{haarLie}(N_{k,r},1)\) along precisely
this map.  Equation (RCV.39) is supported on \(N_{k,r}\), all
of which lies in this image.  Thus it may equivalently be pushed
forward along the inverse conjugation, with no discarded mass.

In the normalized chart, the readings of \(Q\) and \(K\)
are unchanged.  The final entry in each row block is
\(h(f_b,gz)=\sqrt T\,h(f_b,g(z/\sqrt T))\).
There are \(r\) such complex entries.  Multiplying each by
\(\sqrt T\) has real determinant \((\sqrt T)^{2r}=T^r\).
The pushforward of Lebesgue measure under this invertible linear
map is \(T^{-r}\) times Lebesgue measure.  Combining this
with (RCV.39) proves (RCV.41).  The coordinate identities hold
on the supporting closed subgroup; hence they give the required
almost-everywhere measurability for the original ambient reading.

The measure in (RCV.41) is finite on compact subsets, nonzero,
and sigma-finite.  It is supported on the actual unipotent subgroup,
and RCV11 proves left invariance in these coordinates.  The extra
condition \((g-1)z'\in\langle e\rangle\) in the definition of
\(N(\epsilon,e,z')\) follows from \(z'\in e^\perp\).
Therefore this subgroup is the same one for every stated \(z'\).
These observations prove every field of the broad original
\texttt{refNg\_isHaarLoc} statement as well as the full
\texttt{refN\_map\_kapL} formula.
\end{proof}
% END RCV22

% BEGIN RCV23
\paragraph{RCV23. Changing the word: the exact row determinant on all matrices.}
\label{claim:RCV23}
For every \(X\), every pair of words \(u,w\), and every
chosen infinite place \(v\), there is an invertible real linear
map \(L_{u,w,v}\) on the full local coordinate space such that
\[
 \kappa_{w,v}(s)=L_{u,w,v}\kappa_{u,v}(s)
       \quad\hbox{for every ambient unitary element }s,
 \qquad |\det L_{u,w,v}|=(T_w/T_u)^r.
\tag{RCV.42}
\]
The same coordinate identity holds for every ambient matrix on which
the readings are defined by the source's total inverse convention.
In particular this is the full original \texttt{FrameRow}.
\begin{proof}
Both \(S_u=(w_{u,v,1},\ldots,w_{u,v,k},z_u)\) and
\(S_w=(w_{w,v,1},\ldots,w_{w,v,k},z_w)\) are bases of the
same complement \(W_v\).  Write \(S_w=S_u C\).  Their Gram
identities give
\[
       C^*\operatorname{diag}(1_k,-T_u)C
                           =\operatorname{diag}(1_k,-T_w).
\]
Taking determinants proves \(|\det_{\mathbb C}C|^2=T_w/T_u\).
For an arbitrary ambient matrix \(s\), the functional
\(y\mapsto h(f_b,sy)\) is complex linear.  Its row of values
on \(S_w\) is therefore its row of values on \(S_u\), multiplied
on the right by the same matrix \(C\).  This proves the row
transformation without any unipotence assumption on \(s\).

All \(Q\), \(A\), \(M=Q^{-1}A\), and central coordinate
readings use only \(e,f,s\).  They are independent of the word,
even when the total matrix inverse is evaluated off its nonsingular
locus.  Thus \(L_{u,w,v}\) is the identity on those coordinates
and right multiplication by \(C\) on each of the \(r\) row
blocks.  A complex linear map has real determinant
\(|\det_{\mathbb C}C|^2\), as follows for instance by writing
its real block matrix.  The product over the \(r\) rows is
\((T_w/T_u)^r\), proving (RCV.42).
No relation between the separately supplied orthonormal frames is
assumed; their Gram matrices alone determine this determinant.
\end{proof}
% END RCV23

% BEGIN RCV24
\paragraph{RCV24. The complete independent-word reference-box formula.}
\label{claim:RCV24}
For every \(X\) and every independently chosen pair \(u,w\),
let
\[
 \nu_u=\prod_{v\mid\infty}\mathrm{refNX}(X,u,v).
\]
Then \(\nu_u\) is an original \texttt{IsHaarOnSet} measure
on \(N_\infty(w)\), and
\[
       \nu_u(\mathrm{Box}_w)
                 =T_w^{-rd}\,2^{d(rk+r^2)}.
\tag{RCV.43}
\]
This proves the full original \texttt{RefNBox}, including the
reference word different from the lattice word.
\begin{proof}
At each place, RCV22 with the frame of \(u\) gives density
\(2^{r^2+rk}T_u^{-r}\) in the \(u\)-coordinates.  Apply
the linear map in RCV23.  Pushforward of Lebesgue measure divides
the density by its absolute determinant, so in the
\(w\)-coordinates it becomes
\[
 2^{r^2+rk}T_u^{-r}(T_w/T_u)^{-r}
                              =2^{r^2+rk}T_w^{-r}.
\]
This equality keeps the same local reference measure throughout.
The local \(w\)-cube has Lebesgue volume one.  Taking the finite
product over the \(d\) specified places yields (RCV.43), since
the original \(\mathrm{Box}_w\) is exactly the product of these
local boxes.

Each local measure is sigma-finite, nonzero, supported on the local
group, and invariant under its left translations by the final
assertion of RCV22.  Finite products preserve all four properties;
nonzeroness can also be read from the positive finite box mass just
computed.  This proves the full original product Haar assertion
\texttt{refNX\_haar} for every \(u,w\), not merely for equal
words.  Empty row or pair index sets contribute the factors one.
\end{proof}
% END RCV24

% BEGIN RCV25
\paragraph{RCV25. Full reference covolume and the original discriminant return.}
\label{claim:RCV25}
For every original \(X\), every reference word \(u\), every
lattice word \(w\), and every supplied
\(\operatorname{IsSFund}(\mathrm{NOinf}(X,w),\Omega)\),
\[
                  \nu_u(\Omega)=\Delta^{e_N}
                         \quad\hbox{in }(0,\infty).
\tag{RCV.44}
\]
In particular, for every original bridge \(hb\), every two
supplied domains \(\Omega,\Omega'\), and every \(u,u'\)
satisfying the original \(\mathrm{HypX}(X,hb,\Omega,\Omega',u,u')\),
\[
 (\nu_u(\Omega))^{\mathrm{toReal}}
 (\nu_u(\Omega'))^{\mathrm{toReal}}
 =|\operatorname{disc}(X.B.I.F)|^{e_{DF}},\qquad
 e_{DF}=-\frac{8r^3+6r^2+6kr^2+r}{3}.
\tag{RCV.45}
\]
This proves the full original \texttt{covol\_refN} and
\texttt{CovolRef}, with no retained covolume, row-lattice, or
finite-level membership premise.
\begin{proof}
Apply RCV19 to the unchanged measure \(\nu_u\), using its
original Haar fields from RCV24.  Equations (RCV.32), (RCV.35),
and (RCV.43) give
\[
 \begin{split}
 \nu_u(\Omega)
 &=\big(T_w^{rd}2^{-d(rk+r^2)}\Delta^{e_N}\big)
                       \big(T_w^{-rd}2^{d(rk+r^2)}\big)\\
 &=\Delta^{e_N}.
 \end{split}
\]
Every cancelled factor is finite and strictly positive; these
properties follow from \(T_w>0\) and \(\Delta>0\), so there
is no cancellation of zero or infinite extended-real masses.
This proves (RCV.44) for every domain in the broad statement.

The original \(\mathrm{HypX}\) supplies exactly the two required
strict-domain assertions: \(\Omega\) for
\(\mathrm{NOinf}(X,u)\) and \(\Omega'\) for
\(\mathrm{NOinf}(X,u')\).  Apply (RCV.44) twice, with reference
word \(u\) both times, and lattice words \(u\) and \(u'\)
respectively.  Finite positive masses have their ordinary real values
under \(\mathrm{toReal}\).  Their product is
\(\Delta^{2e_N}=\Delta^{e_{DF}}\), giving (RCV.45).
The bridge and the other fields of \(\mathrm{HypX}\) are accepted
at their original scope; this calculation does not need to constrain
or reselect them.

The proof has thus supplied the exact reference volume that the
source's subsequent \texttt{placeSplit\_of} consumes.  That next
theorem still has its separate local \texttt{JacGauge} input.
The present result does not assert its Jacobian or the full
\texttt{RealBracket}, and it does not reconstruct \texttt{Forms0}
or any missing global spectral producer.
\end{proof}
% END RCV25

\subsection{What has been fixed for the next comparison}

The completed formulas have distinct purposes.  Equation (RCV.35)
works with every supplied Haar normalization and every original
arithmetic domain.  Equation (RCV.43) computes the fixed reference
normalization, with its reference word independent of the word in
the box.  Their cancellation gives the discriminant in (RCV.45).
This distinction is necessary in the next step: local disintegration
must compare the actual measures with these reference measures,
while keeping the same setting, two words, domains, and bridge.

The remaining analytic card starts at the original
\texttt{JacGauge} and its all-section local disintegration formula.
It may then use the proved \texttt{CovolRef} in
\texttt{placeSplit\_of}, and the independently proved or separately
reconstructed product-disintegration supplier in the exact
\texttt{RealBracket} consumer.  Neither a change of an infinity
measure nor the already completed actual \texttt{InVolumes}
can substitute for that local Jacobian calculation.
